\documentclass[a4paper,12pt]{article}

% -------------------------------------------------
% Кодировка и язык
% -------------------------------------------------
\usepackage[T2A]{fontenc}
\usepackage[utf8]{inputenc}
\usepackage[russian]{babel}

% -------------------------------------------------
% Математика
% -------------------------------------------------
\usepackage{amsmath,amssymb,mathtools}
\usepackage{icomma}

% -------------------------------------------------
% Типографика
% -------------------------------------------------
\usepackage{helvet}
\renewcommand{\familydefault}{\sfdefault}
\usepackage{microtype}
\usepackage{geometry}
\geometry{
  left=20mm,
  right=20mm,
  top=20mm,
  bottom=22mm
}
\usepackage{parskip}
\usepackage{setspace}
\onehalfspacing

% -------------------------------------------------
% Таблицы и списки
% -------------------------------------------------
\usepackage{tabularx,booktabs}
\usepackage{enumitem}

% -------------------------------------------------
% Графика
% -------------------------------------------------
\usepackage{tikz}
\usetikzlibrary{calc}
\usetikzlibrary{arrows.meta,positioning,shapes.geometric}

% -------------------------------------------------
% Служебное
% -------------------------------------------------
\usepackage{xcolor}
\usepackage{fancyhdr}
\usepackage{hyperref}

% -------------------------------------------------
% Цветовая палитра
% -------------------------------------------------
\definecolor{accent}{HTML}{008080}
\definecolor{darkaccent}{HTML}{004D4D}
\definecolor{textgray}{HTML}{333333}
\definecolor{softgray}{HTML}{747474}
\definecolor{linegray}{HTML}{D9E1E1}

\color{textgray}

% -------------------------------------------------
% PDF
% -------------------------------------------------
\hypersetup{
  colorlinks=true,
  linkcolor=darkaccent,
  urlcolor=darkaccent,
  pdftitle={Чек-лист. Тригонометрические формулы и уравнения},
  pdfauthor={Лёвин Артём Александрович}
}

% -------------------------------------------------
% Колонтитулы
% -------------------------------------------------
\pagestyle{fancy}
\fancyhf{}
\fancyhead[L]{\small Тригонометрия}
\fancyhead[R]{\small Ярослав Гаврилов}
\fancyfoot[C]{\small\thepage}
\renewcommand{\headrulewidth}{0.3pt}
\renewcommand{\footrulewidth}{0pt}

% -------------------------------------------------
% Заголовки
% -------------------------------------------------
\newcommand{\maintitle}[1]{%
  {\Large\bfseries\color{darkaccent}#1\par}
  \vspace{0.5em}
  {\color{accent}\rule{\linewidth}{0.7pt}\par}
  \vspace{0.8em}
}

\newcommand{\subheading}[1]{%
  \vspace{0.7em}
  {\large\bfseries\color{darkaccent}#1\par}
  \vspace{0.25em}
}

\newcommand{\important}[1]{%
  \par\noindent
  {\bfseries\color{darkaccent}#1}
}

% -------------------------------------------------
% TikZ: блок-схемы
% -------------------------------------------------
\tikzset{
  flowstart/.style={
    ellipse,
    draw=accent,
    line width=0.9pt,
    minimum width=42mm,
    minimum height=10mm,
    align=center,
    font=\small
  },
  flowaction/.style={
    rectangle,
    rounded corners=2mm,
    draw=accent,
    line width=0.8pt,
    text width=68mm,
    minimum height=12mm,
    align=center,
    inner sep=4pt,
    font=\small
  },
  flowdecision/.style={
    diamond,
    draw=darkaccent,
    line width=0.8pt,
    aspect=2.5,
    text width=43mm,
    align=center,
    inner sep=2pt,
    font=\small
  },
  flowarrow/.style={
    -{Stealth[length=3mm,width=2mm]},
    line width=0.8pt,
    draw=darkaccent
  },
  flowlabel/.style={
    font=\small,
    fill=white,
    inner sep=1.5pt
  }
}

\setlist[itemize]{
  leftmargin=1.7em,
  itemsep=0.35em,
  topsep=0.35em
}

\setlist[enumerate]{
  leftmargin=2em,
  itemsep=0.55em,
  topsep=0.45em
}

\begin{document}

% =================================================
% ТИТУЛЬНЫЙ ЛИСТ
% =================================================
\begin{titlepage}
\thispagestyle{empty}

\begin{center}

\vspace*{24mm}

{\color{accent}\rule{32mm}{1.2pt}}

\vspace{11mm}

{\Huge\bfseries\color{darkaccent}
Чек-лист
\par}

\vspace{7mm}

{\LARGE\bfseries
Тригонометрические формулы\\[2mm]
и решение тригонометрических уравнений
\par}

\vspace{12mm}

{\large
Основные тождества, формулы суммы и разности,\\
формулы двойного угла, приведение к одной функции,\\
замена и решение уравнений
\par}

\vfill

{\large
Ярослав Гаврилов
\par}

\vspace{7mm}

{\normalsize
Подготовка к ЕГЭ по профильной математике
\par}

\vspace{12mm}

{\normalsize
\today
\par}

\vspace{18mm}

{\small
Лёвин Артём Александрович\\
эксклюзивно для Ярослава Гаврилова
\par}

\vspace{21mm}

\end{center}

% Сдержанная кривая Безье внизу титульного листа
\begin{tikzpicture}[remember picture,overlay]
  \draw[
    accent!55,
    line width=1.4pt
  ]
  ($(current page.south west)+(0mm,15mm)$)
  .. controls
  ($(current page.south west)+(55mm,25mm)$)
  and
  ($(current page.south east)+(-60mm,7mm)$)
  ..
  ($(current page.south east)+(0mm,17mm)$);
\end{tikzpicture}

\end{titlepage}

% =================================================
% ОСНОВНОЙ ЧЕК-ЛИСТ
% =================================================
\maintitle{1. Основное тригонометрическое тождество}

Главная формула:

\[
\boxed{\sin^2\alpha+\cos^2\alpha=1}
\]

Она верна для любого допустимого значения угла \(\alpha\).

\subheading{Что необходимо увидеть}

Чтобы непосредственно получить единицу, должны одновременно выполняться условия:

\begin{itemize}
  \item аргументы синуса и косинуса совпадают;
  \item синус и косинус возведены во вторую степень;
  \item коэффициенты перед \(\sin^2\alpha\) и \(\cos^2\alpha\) равны \(1\).
\end{itemize}

Например,

\[
\sin^2 37^\circ+\cos^2 37^\circ=1.
\]

Значение самого угла здесь роли не играет.

При разных аргументах такое преобразование недопустимо:

\[
\sin^2 60^\circ+\cos^2 59^\circ\ne 1
\quad\text{в общем случае.}
\]

\subheading{Если коэффициенты мешают}

Иногда тождество можно сначала \emph{выделить}.

Например,

\[
4\sin^2\alpha+\cos^2\alpha
=
3\sin^2\alpha+
\bigl(\sin^2\alpha+\cos^2\alpha\bigr)
=
3\sin^2\alpha+1.
\]

Если коэффициенты одинаковы, общий множитель удобно вынести:

\[
4\sin^2\alpha+4\cos^2\alpha
=
4\bigl(\sin^2\alpha+\cos^2\alpha\bigr)
=
4.
\]

\subheading{Формулу нужно знать в обе стороны}

Полезно видеть не только

\[
\sin^2\alpha+\cos^2\alpha\longrightarrow 1,
\]

но и обратное преобразование:

\[
1\longrightarrow \sin^2\alpha+\cos^2\alpha.
\]

Поэтому любое число \(A\) можно при необходимости представить как

\[
A
=
A\bigl(\sin^2\alpha+\cos^2\alpha\bigr).
\]

Например,

\[
5
=
5\sin^2 20^\circ+5\cos^2 20^\circ.
\]

И аналогично

\[
-5
=
-5\sin^2 120^\circ-5\cos^2 120^\circ.
\]

\important{Ключевая идея.}
Тригонометрические формулы полезно узнавать в обоих направлениях:
раскрывать их и, наоборот, сворачивать.

% =================================================
\maintitle{2. Формулы суммы и разности}

Для синуса:

\[
\boxed{
\sin(\alpha+\beta)
=
\sin\alpha\cos\beta+
\cos\alpha\sin\beta
}
\]

\[
\boxed{
\sin(\alpha-\beta)
=
\sin\alpha\cos\beta-
\cos\alpha\sin\beta
}
\]

Для косинуса:

\[
\boxed{
\cos(\alpha+\beta)
=
\cos\alpha\cos\beta-
\sin\alpha\sin\beta
}
\]

\[
\boxed{
\cos(\alpha-\beta)
=
\cos\alpha\cos\beta+
\sin\alpha\sin\beta
}
\]

\subheading{Зачем они нужны}

Формулы позволяют:

\begin{itemize}
  \item вычислять нетабличные значения через табличные углы;
  \item раскрывать сложные аргументы;
  \item преобразовывать выражения так, чтобы аргументы в уравнении стали одинаковыми.
\end{itemize}

\subheading{Пример: \(\sin 15^\circ\)}

Представим

\[
15^\circ=45^\circ-30^\circ.
\]

Тогда

\[
\begin{aligned}
\sin 15^\circ
&=
\sin(45^\circ-30^\circ)
\\
&=
\sin45^\circ\cos30^\circ
-
\cos45^\circ\sin30^\circ
\\
&=
\frac{\sqrt2}{2}\cdot\frac{\sqrt3}{2}
-
\frac{\sqrt2}{2}\cdot\frac12
\\
&=
\frac{\sqrt6-\sqrt2}{4}.
\end{aligned}
\]

Аналогично

\[
\sin75^\circ
=
\sin(45^\circ+30^\circ)
=
\frac{\sqrt6+\sqrt2}{4}.
\]

\subheading{Частный полезный случай}

Рассмотрим

\[
\sin\left(\frac{\pi}{2}+x\right).
\]

По формуле суммы:

\[
\begin{aligned}
\sin\left(\frac{\pi}{2}+x\right)
&=
\sin\frac{\pi}{2}\cos x+
\cos\frac{\pi}{2}\sin x
\\
&=
1\cdot\cos x+0\cdot\sin x
\\
&=
\cos x.
\end{aligned}
\]

Следовательно,

\[
\boxed{
\sin\left(\frac{\pi}{2}+x\right)=\cos x
}
\]

Аналогично,

\[
\boxed{
\cos\left(\frac{\pi}{2}-x\right)=\sin x
}
\]

Эти равенства можно получать через формулы приведения, однако
формулы суммы и разности дают надёжный универсальный способ проверки.

% =================================================
\maintitle{3. Формулы двойного угла}

\subheading{Синус двойного угла}

\[
\boxed{
\sin 2\alpha=2\sin\alpha\cos\alpha
}
\]

Формула следует из формулы суммы:

\[
\begin{aligned}
\sin2\alpha
&=
\sin(\alpha+\alpha)
\\
&=
\sin\alpha\cos\alpha+
\cos\alpha\sin\alpha
\\
&=
2\sin\alpha\cos\alpha.
\end{aligned}
\]

\subheading{Главная практическая идея}

Если перед Вами

\[
\sin u,
\]

то формулу двойного угла можно записать как

\[
\boxed{
\sin u
=
2\sin\frac{u}{2}\cos\frac{u}{2}
}
\]

То есть исходный аргумент целиком делится на \(2\).

Например,

\[
\sin3\alpha
=
2\sin\frac{3\alpha}{2}
\cos\frac{3\alpha}{2},
\]

\[
\sin101\alpha
=
2\sin\frac{101\alpha}{2}
\cos\frac{101\alpha}{2}.
\]

Формула работает и в обратную сторону:

\[
2\sin3\alpha\cos3\alpha
=
\sin6\alpha.
\]

\subheading{Если множителя \(2\) нет}

Выражение

\[
\sin5\alpha\cos5\alpha
\]

ещё нельзя непосредственно свернуть по формуле двойного угла.

Добавим нужную двойку:

\[
\sin5\alpha\cos5\alpha
=
\frac12
\bigl(
2\sin5\alpha\cos5\alpha
\bigr).
\]

Следовательно,

\[
\boxed{
\sin5\alpha\cos5\alpha
=
\frac12\sin10\alpha
}
\]

В общем случае

\[
\boxed{
\sin u\cos u=\frac12\sin2u
}
\]

\subheading{Косинус двойного угла}

Необходимо знать три равносильные формы:

\[
\boxed{
\cos2\alpha
=
\cos^2\alpha-\sin^2\alpha
}
\]

\[
\boxed{
\cos2\alpha
=
2\cos^2\alpha-1
}
\]

\[
\boxed{
\cos2\alpha
=
1-2\sin^2\alpha
}
\]

Для произвольного аргумента \(u\):

\[
\cos u
=
\cos^2\frac{u}{2}
-
\sin^2\frac{u}{2},
\]

\[
\cos u
=
2\cos^2\frac{u}{2}-1,
\]

\[
\cos u
=
1-2\sin^2\frac{u}{2}.
\]

Например,

\[
\cos3\alpha
=
2\cos^2\frac{3\alpha}{2}-1.
\]

\important{Как выбрать форму?}
Смотрите на остальные части задачи.
Если выгодно получить только косинусы, используйте
\(2\cos^2\alpha-1\).
Если нужны только синусы, используйте
\(1-2\sin^2\alpha\).

% =================================================
\maintitle{4. Главная концепция решения тригонометрических уравнений}

Перед алгебраическим решением удобно стремиться к двум целям:

\[
\boxed{\text{одинаковые аргументы}}
\]

и

\[
\boxed{\text{одинаковая тригонометрическая функция}}
\]

Если в одной части уравнения присутствует \(2x\), а в другой \(x\),
формулы двойного угла часто позволяют перейти от \(2x\) к \(x\).

Если одновременно присутствуют синусы и косинусы,
основное тригонометрическое тождество и формулы двойного угла
часто позволяют оставить только одну функцию.

\subheading{Опорный пример}

Решим

\[
\cos2x
=
\sin\left(x+\frac{\pi}{2}\right).
\]

Сначала преобразуем правую часть:

\[
\sin\left(x+\frac{\pi}{2}\right)=\cos x.
\]

Получаем

\[
\cos2x=\cos x.
\]

Справа уже стоит косинус, поэтому для левой части рационально выбрать

\[
\cos2x=2\cos^2x-1.
\]

Тогда

\[
2\cos^2x-1=\cos x.
\]

Переносим всё в одну часть:

\[
2\cos^2x-\cos x-1=0.
\]

Теперь уравнение стало алгебраическим относительно \(\cos x\).

Выполняем замену

\[
t=\cos x,
\qquad
-1\le t\le1.
\]

Получаем

\[
2t^2-t-1=0.
\]

Корни:

\[
t_1=1,
\qquad
t_2=-\frac12.
\]

Оба удовлетворяют ограничению

\[
-1\le t\le1.
\]

Возвращаемся к исходной переменной:

\[
\cos x=1
\]

или

\[
\cos x=-\frac12.
\]

Следовательно,

\[
\boxed{
x=2\pi n,\qquad n\in\mathbb Z
}
\]

или

\[
\boxed{
x=\pm\frac{2\pi}{3}+2\pi n,
\qquad n\in\mathbb Z
}
\]

\subheading{Почему ограничение для замены обязательно}

При

\[
t=\sin x
\quad\text{или}\quad
t=\cos x
\]

всегда выполняется

\[
\boxed{-1\le t\le1}.
\]

Поэтому, например, корень

\[
t=2
\]

алгебраического уравнения необходимо отбросить:
значения \(\sin x=2\) и \(\cos x=2\) невозможны.

\subheading{Ограничения и типичные риски}

\begin{itemize}
  \item До преобразований проверяйте ОДЗ, если в уравнении есть дроби, корни или другие выражения с ограничениями.
  \item Не делите уравнение на \(\sin x\), \(\cos x\) или другое выражение с переменной без отдельного рассмотрения случая, когда делитель равен нулю.
  \item При замене \(t=\sin x\) или \(t=\cos x\) обязательно учитывайте \(t\in[-1;1]\).
  \item После возведения обеих частей уравнения в квадрат обязательна проверка найденных корней.
  \item Следите за знаками при переносе слагаемых.
  \item Формула применяется только при выполнении её структуры: аргументы и коэффициенты должны соответствовать формуле.
\end{itemize}

\important{Рабочий принцип.}
Если преобразование делает аргументы и функции более однородными,
обычно Вы движетесь в полезном направлении.


% =================================================
% КРАТКИЙ ЧЕК-ЛИСТ
% =================================================
\maintitle{5. Чек-лист перед решением}

Перед тем как переходить к вычислениям, проверьте:

\begin{itemize}
  \item[$\square$] Я определил ОДЗ и ограничения, если они есть.
  \item[$\square$] Я сравнил аргументы всех тригонометрических функций.
  \item[$\square$] Я понимаю, какая формула позволяет сделать аргументы одинаковыми.
  \item[$\square$] Я стараюсь оставить одну основную функцию: синус или косинус.
  \item[$\square$] Я узнаю формулы как слева направо, так и справа налево.
  \item[$\square$] Я перенёс всё в одну часть, если это помогает увидеть структуру.
  \item[$\square$] При замене \(t=\sin x\) или \(t=\cos x\) я записал \(t\in[-1;1]\).
  \item[$\square$] Я не делил на выражение, которое может быть равно нулю, без проверки отдельного случая.
  \item[$\square$] После обратной замены я записал все серии решений.
\end{itemize}

\subheading{Минимум формул, который необходимо узнавать сразу}

\[
\sin^2x+\cos^2x=1
\]

\[
\sin(\alpha\pm\beta)
=
\sin\alpha\cos\beta
\pm
\cos\alpha\sin\beta
\]

\[
\cos(\alpha\pm\beta)
=
\cos\alpha\cos\beta
\mp
\sin\alpha\sin\beta
\]

\[
\sin2x=2\sin x\cos x
\]

\[
\cos2x
=
\cos^2x-\sin^2x
=
2\cos^2x-1
=
1-2\sin^2x.
\]

% =================================================
% ТРЕНИРОВОЧНЫЙ БЛОК
% =================================================
\clearpage
\thispagestyle{fancy}

\maintitle{Тренировочный блок. Путь А}

\textbf{Задача 1. Базовые преобразования}

Упростите выражения:

\begin{enumerate}[label=\alph*)]
  \item
  \[
  4\sin^2\alpha+4\cos^2\alpha;
  \]

  \item
  \[
  2\sin5x\cos5x;
  \]

  \item представьте
  \[
  \cos3x
  \]
  только через \(\cos\dfrac{3x}{2}\).
\end{enumerate}

\vspace{8mm}

\textbf{Задача 2. Формула суммы}

Вычислите точное значение

\[
\sin75^\circ.
\]

Ответ представьте в виде выражения с корнями.

\vspace{10mm}

\textbf{Задача 3. Тригонометрическое уравнение}

Решите уравнение

\[
\cos2x+\sin\left(x+\frac{\pi}{2}\right)=0.
\]

Запишите все решения.

% =================================================
% ОТВЕТЫ
% =================================================
\clearpage
\thispagestyle{fancy}

\maintitle{Ответы}

\renewcommand{\arraystretch}{1.45}

\begin{table}[ht]
\centering
\begin{tabularx}{\linewidth}{
  >{\raggedright\arraybackslash}p{22mm}
  X
}
\toprule
\textbf{Задание} & \textbf{Ответ} \\
\midrule

1а &
\[
4
\]
\\

1б &
\[
\sin10x
\]
\\

1в &
\[
2\cos^2\frac{3x}{2}-1
\]
\\

2 &
\[
\frac{\sqrt6+\sqrt2}{4}
\]
\\

3 &
\[
x=\pm\frac{\pi}{3}+2\pi n
\quad\text{или}\quad
x=\pi+2\pi n,
\qquad n\in\mathbb Z
\]
\\

\bottomrule
\end{tabularx}
\end{table}

\vfill

\begin{center}
{\small\color{softgray}
Лёвин Артём Александрович
}
\end{center}

\end{document}